\documentclass{article}

% Required packages for drawing and animation
\usepackage{tikz}
\usepackage{animate}

% Set page geometry for a better layout
\usepackage[a4paper, margin=1in]{geometry}

\begin{document}

\begin{center}
\section*{Animation of an Archimedean Spiral Construction}

% Begin the animation environment.
% 'autoplay' makes it start automatically.
% 'loop' makes it repeat.
% '25' is the frame rate (frames per second).
\begin{animateinline}[autoplay, loop]{25}
    % This loop iterates through the angles from 0 to 1080 degrees in steps of 5.
    % Each iteration will create a new frame in the animation.
    \foreach \angle in {0, 5, ..., 1080}{
        % Create a new frame for the animation
        \newframe
        
        % Start the TikZ picture environment for drawing
        \begin{tikzpicture}
            % Set the boundary of the drawing area.
            % The values are determined by the maximum radius of the spiral.
            % 1080 degrees * 0.005 = 5.4 radians.
            \useasboundingbox (-6,-6) rectangle (6,6);

            % 1. Draw the spiral path up to the current angle
            % The plot is drawn in polar coordinates (angle:radius).
            % The radius '0.005*\a' increases linearly with the angle '\a',
            % which is the definition of an Archimedean spiral (r = aθ).
            \draw[cyan, thick, domain=0:\angle, variable=\a, samples=200]
                plot (\a: {0.005*\a});

            % 2. Draw the "string" (a rotating radius line)
            % This is a dashed red line from the origin (0,0) to the current
            % point on the spiral.
            \draw[red, dashed] (0,0) -- (\angle:{0.005*\angle});

            % 3. Draw the "pen" (a point at the end of the string)
            % This is a filled red circle at the current point on the spiral.
            \fill[red] (\angle:{0.005*\angle}) circle (2pt);
            
            % 4. (Optional) Draw coordinate axes for reference
            \draw[->, gray, thin] (-6,0) -- (6,0) node[right] {$x$};
            \draw[->, gray, thin] (0,-6) -- (0,6) node[above] {$y$};
        \end{tikzpicture}
    }
\end{animateinline}

\subsection*{How it Works}
This animation is created by drawing a sequence of frames. In each frame:
\begin{enumerate}
    \item A small segment of the \textcolor{cyan}{cyan spiral} is drawn. The radius of the spiral is directly proportional to the angle, which defines the Archimedean spiral.
    \item A \textcolor{red}{dashed red line} represents the "string" unwinding from the center to the current point.
    \item A \textcolor{red}{red dot} represents the pen at the tip of the string.
\end{enumerate}
To see the animation, open the compiled PDF in Adobe Acrobat Reader.

\end{center}

\end{document}
